\section{A finite comparison for the first gap}\label{sec:finite}
\subsection{A uniform bound from finite data at one time}
Take the first positive edge, \(X=4\), \(s=1000-t\in[0,1]\), and \(b=671/500\). Write \(P=U-b,\ W=U^2-b^2\), and let \(\E_0\) here denote expectation at the upper anchor \(1000\). Tilting this one source gives the moments
\begin{equation}
 d_j(s)=\frac{\E_0[P^je^{-sW}]}{\E_0[e^{-sW}]},\quad d_0=1,\quad
 \mu=b+d_1,\quad \sigma^2=d_2-d_1^2. \label{eq:tilt}
\end{equation}
The numerator for each \(0\le j\le6\) is enclosed by a cubic polynomial with absolute remainder
\begin{equation}
 \left|\E_0[P^je^{-sW}]
    -\sum_{m=0}^3\frac{(-s)^m}{m!}\E_0[P^jW^m]\right|
 \le\frac{s^4e^{b^2}}{24}\sqrt{\E_0[P^{2j}W^8]} . \label{eq:tiltrest}
\end{equation}
Indeed \((-sW)_+\le b^2\), and Taylor's integral remainder followed by Cauchy--Schwarz gives \eqref{eq:tiltrest}. Since \(W=P(2b+P)\), all source integrals are anchor moments through degree 28. The two-copy law gives, for \(p=2,4\),
\begin{equation}
 m_p^*=\frac1\mu\sum_{i=0}^p(-1)^i\binom pi d_i
                    \{b\,d_{p-i}+d_{p-i+1}\}. \label{eq:certmom}
\end{equation}
This expands \(\E[U(U-U')^p]/\mu\), using the exchange symmetry for even \(p\).

All moments were evaluated by outward Arb arithmetic at 224 bits, using the original anchor integrator and its complete theta tail allowance. It integrates the first square mode on \([1,2]\), bounds the higher square modes there by geometric comparison, and bounds the two omitted half-line intervals by the saddle estimates in \cite[\S6]{Founding08}. Its absolute allowance is below \(2.353\,10^{-107}\) in each peak-scaled integral for functions bounded by \((u+2)^7e^{16(u+2)}\). This covers \(|P|^{28}\) and all lower powers, since \(21\log(u+2)\le (21/e)(u+2)\). It also covers the real trigonometric endpoint tests below.

Evaluate \eqref{eq:tiltrest} on 16 equal closed subintervals of \(s\). Each tilted mass has a positive lower bound. On every subinterval \(j(t,x)>0.99817\) throughout \([0,4]\). Also \(\left|\E e^{\ii x(U-\mu)}-1\right|\le x^2\sigma^2/2\); the checked inequality \(\mu X-\arcsin(X^2\sigma^2/2)>3\pi/2\) places both labelled zeros inside the interval. The centered transform remains in the corresponding disk around one, which fixes the argument from zero. These are whole-cell inequalities.

\Needspace{9\baselineskip}
\begin{proposition}[Finite signed descent]\label{prop:check}
For \(999\le t\le1000\) the first positive edge satisfies
\begin{equation}
 0.0002270493<-\dot g(t)/g(t)<0.0002300072.
\end{equation}
Its anchor value is \(D(1000)=1.3709993248674981982\ldots\).
Integrating the cell-specific bounds in Proposition~\ref{prop:finite} gives
\begin{equation}
 1.37162235<D(999)<1.37162972. \label{eq:number}
\end{equation}
\end{proposition}
\begin{proof}
Apply \eqref{eq:certmom} on each cell, check \(b_*>0,l>0\), and multiply the anchor \(D\) by the interval exponentials in \eqref{eq:slab}. The cell widths are exact dyadics. The program, precision, sign tests and outward endpoints are recorded alongside the manuscript \cite{FiniteRecord}.
\end{proof}
An equally informed endpoint calculation uses the same \(d_0,\ldots,d_6\) at \(s=1\). At real \(x\) evaluate
\[
 \sum_{j=0}^5\frac{x^jd_j}{j!}\cos(bx+j\pi/2)
 \quad\hbox{with error at most } |x|^6 d_6/720 .
\]
Signs on opposite sides and the common phase count enclose each root. This gives \(1.3716257900<D(999)<1.3716262696\), width below \(4.80\,10^{-7}\), versus \(7.36\,10^{-6}\) for transport. Direct quadrature of the complete upper-anchor field with the multiplier \(e^{b^2-U^2}\), using separate endpoint signs and the same count, gives \(D(999)=1.37162603031689388115504\ldots\). It shares source and tail code with the moment calculation. The path estimate establishes strict widening at every intermediate time. For the endpoint alone the finite Taylor comparison is tighter.

\subsection{Removing the change in mean spacing}
On a counted chart let \(h=\log j\) and let angle brackets denote spatial averaging along the edge. From \eqref{eq:difference} and evenness,
\begin{equation}
 \frac{\dd}{\dd t}\log(a g)
 =\left\langle\int_0^x(x-y)h^{(4)}(t,y)\dd y\right\rangle,\qquad
 h^{(4)}(t,0)=m_4^*-3w^2 . \label{eq:quartic}
\end{equation}
The superscript denotes spatial differentiation. Freezing the last derivative at zero gives the explicit edge term
\((m_4^*-3w^2)(x_-^2+x_-x_++x_+^2)/6\).
If \(\sup_{[0,X]}|h^{(6)}|\le B_6\), evenness and Taylor's theorem bound the error by \(B_6 X^4/24\). Thus a shared signed fourth cumulant and a controlled sixth derivative can replace separate absolute payments for \(j''\) and \((j')^2\). A later calculation with the same 28 upper-time moments carried this comparison to \(t=984\) for the first edge. Direct evaluation gave a tighter endpoint interval. Both methods certified widening along this already large-time path \cite{EZJoined}.

The material acceleration on this chart satisfies \eqref{eq:acceleration}. A bound on \(K\) together with initial \(D,V\) enters the parent's stopped-path barrier \cite[\S4]{AGC}. The divided differences \eqref{eq:Sgen}--\eqref{eq:Tgen} give a separate way of comparing that forcing with fixed-order exterior moments. Through a loss of the phase chart the parent's count-two contour \cite[\S\S2,6]{AGC} still requires a complete boundary floor and a paid error for its center and discriminant. Relabelling a new pair at the end of a step would invalidate the comparison.

A simple control tests these hypotheses. For \(H(t,x)=e^t\cos x+c\), \(c>1\), a positive half-line lift is \(B=e^{t+\ii x}+c\). Its current is \(e^{2t}+ce^t\cos x\). At \(ce^{-t}=1/2\), the first two roots are \(2\pi/3,4\pi/3\), and their physical gap velocity is \(2/\sqrt3>0\). The current is positive there, but \(m_4^*=w=1/3\); any \(X\ge4\pi/3\) fails \(w-X^2m_4^*/2>0\). At \(t=\log c,\ x=\pi\) the current vanishes as the two roots meet. The single-cosine control has constant positive current, \(w=0\), and stationary gaps. Both satisfy the representation identities. The signed theta estimate uses quantitative source concentration. Weighted-square and finite-prime controls change the normalized theta weights, reflected smooth controls lack the integer-square expansion, and altered observations change the cosine readout. The source and count assumptions for the additional controls are detailed in \cite{EZPhaseScope,EZMixedAnchor}. The finite computation is inside a previously all-real regime. Lower-time and all-edge coverage remain open.
